\scanpage{463}\mspage{375}
% archive sheet number: 462
\deleted{\unsure{$X$}}
\[
\underline{\mathrm{Hom}}_{\Phi,\pi_0}^{\Psi}(M,N)\xrightarrow{\approx}
\underline{\mathrm{Hom}}_{\Phi'}^{\Psi}(C,N),
\]
où $\Phi'$ est \unsure{formée|formé} des catégories qui sont \deleted{soit dans} $\Phi$, ou qui sont gdes devant $\pi_0$ et de cardinal $\leq\pi$. Il suffit \unsure{donc} \deleted{\unsure{de}}d'appliquer encore \unsure{4.14.2|4.14.?}.

\uline{Corollaire} 4.14.5. \deleted{Soit $\Phi$ l'un des} \inserted{Considérons les} trois ensembles suivants : a) $\Phi_1$ est formée de toutes les catégories \unsure{*} équivalentes à des petites catégories \quad b) \deleted{On suppose \unsure{toujours}} \inserted{$\Phi_2$ comme a), mais en} se restreignant aux catégories filtrantes.

% slanted
\margin{On pourrait \deleted{\unsure{en fait|en tout}} faire ici un énoncé qui englobe 4.14.4, en rajoutant un d) $\Phi_4$ formé des catégories gdes devant $\pi_0$ et prendre $\Phi$ \deleted{de la forme \unsure{$\Phi$} \unsure{de}} \inserted{\unsure{de la forme|l'un des}} $\Phi_1,\deleted{\unsure{\Phi_2}}\Phi_2,\Phi_4,\Phi_3,$ \deleted{\unsure{ou former}} $\Phi_3\cup\Phi_2,\Phi_3\cup\Phi_4$ \deleted{\unsure{avec ou} $\Phi_2$ \unsure{ou} $\Phi_4$ ou $\Phi_3,\Phi_2$} ou bien la réunion d'un de ceux-ci et d'un petit ensemble $\Phi_0$ \unsure{de|des} cat. \unsure{ess. petites|ens. petites|$\mathcal U$-petites}.}

c) $\Phi_3$ \deleted{est} formé des catégories discrètes, \unsure{de|ou} celles qui sont gdes devant $\pi_0$ ($\pi_0$ un cardinal infini donné). Soit $\Phi$ \deleted{un \unsure{ensemble} de catégories, \unsure{soumises} \unsure{ici}} \inserted{l'un des ens. de catégories} $\Phi_1,\Phi_2,\Phi_3,\Phi_2\cup\Phi_3$, \unsure{;} soit $\Psi$ un \inserted{petit} ensemble de catégories. Soient $M,N$ \deleted{deux} \unsure{catégorie|catégories} \unsure{accessibles}, stables par $\underset{\deleted{\unsure{\gets}}\to}{\lim}$ de type $\Phi$, et par $\varprojlim$ de type $\Psi$. Alors \deleted{pour tte catégorie accessible $N$,} $\underline{\mathrm{Hom}}_\Phi^\Psi(M,N)$ est accessible.
